\documentclass{amsart}
\usepackage{amsmath,amssymb,amsthm,hyperref}

\newtheorem{theorem}{Theorem}
\newtheorem{lemma}{Lemma}
\newtheorem{proposition}{Proposition}
\newtheorem{corollary}{Corollary}
\theoremstyle{definition}
\newtheorem{remark}{Remark}

\DeclareMathOperator{\SL}{SL}
\DeclareMathOperator{\GL}{GL}
\DeclareMathOperator{\Spec}{Spec}
\DeclareMathOperator{\Proj}{Proj}
\DeclareMathOperator{\Sym}{Sym}
\DeclareMathOperator{\disc}{disc}
\DeclareMathOperator{\im}{im}
\newcommand{\PP}{\mathbb{P}}
\newcommand{\CC}{\mathbb{C}}
\newcommand{\nab}{\nabla}

\begin{document}

\title{Normality of the image of the projective gradient map}
\maketitle

\section{Notation and statement}

Throughout, $n\ge 2$ and $d\ge 3$ are integers and we work over $\CC$.
Let $V=\CC^n$ with coordinates $x_1,\dots,x_n$ and $A=\CC[x_1,\dots,x_n]$; write
$A_k=S^k(V)^{*}$ for the space of homogeneous forms of degree $k$ on $V$, a finite
dimensional $\CC$--vector space. We use:
\begin{itemize}
\item $\mathcal V:=A_d$ and $\mathcal W:=(A_{d-1})^{\,n}$, the two ambient vector spaces.
\item $X=X_n^d\subset \mathcal V$, the open set of forms with $\disc(f)\neq 0$ (smooth
hypersurfaces in $\PP^{n-1}$).
\item $Z=Z_n^{d-1}\subset \mathcal W$, the open set of $n$-tuples forming a regular sequence.
\item $\nab\colon \mathcal V\to \mathcal W$, $\nab f=(\partial_1 f,\dots,\partial_n f)$,
$\partial_i=\partial/\partial x_i$. This is a \emph{linear} map of vector spaces; it carries
$X$ into $Z$.
\item $\GL_n=\GL(V)$ acting by $(g\cdot h)(v)=h(g^{-1}v)$ on each $A_k$, and on $\mathcal W$
by \eqref{eq:Zaction} below. We write $G:=\SL_n$.
\end{itemize}

Projectivizing and taking the GIT quotient by $G$ produces, by hypothesis, a finite
injective morphism
\[
\PP\nab\colon M:=\PP X/G \longrightarrow N:=\PP Z/\!/G .
\]
Because $\PP\nab$ is finite, it is proper, so its image is closed; let
$Y:=\PP\nab(M)\subset N$ be this closed image. We prove:

\begin{theorem}\label{thm:main}
For all $n\ge 2$, $d\ge 3$ the image $Y=\PP\nab(M)$ is a normal subvariety of $N$.
\end{theorem}

The proof has three ingredients:
\begin{enumerate}
\item[(A)] The source $M$ is a normal variety (Section~\ref{sec:source}).
\item[(B)] $\PP\nab$ is an isomorphism onto $Y$ (Sections~\ref{sec:euler}--\ref{sec:proj}).
This is the main content; we prove it by an invariant-theoretic argument resting on the fact
that $\nab$ is a \emph{split} linear injection (it has a linear left inverse given by Euler's
identity) together with linear reductivity of $G$.
\item[(C)] Combining (A) and (B): an isomorphism onto a closed image with normal source has
normal image. We also note the converse via Zariski's Main Theorem, recovering the remark in
the problem.
\end{enumerate}

\section{The source is a normal variety}\label{sec:source}

\begin{proposition}\label{prop:sourcenormal}
$M=\PP X/G$ is a normal quasi-projective variety and $\PP X\to M$ is a geometric quotient.
\end{proposition}

\begin{proof}
\emph{Finite stabilizers.} If $g\in G$ fixes $[f]\in \PP X$ then $g\cdot f=\lambda f$ for
some $\lambda\in\CC^{*}$, so $g$ preserves the smooth hypersurface $\{f=0\}\subset\PP^{n-1}$.
By the theorem of Matsumura--Monsky \cite{MM}, the group of linear automorphisms of
$\PP^{n-1}$ preserving a smooth hypersurface of degree $d\ge 3$ is finite for all $n\ge 2$.
Hence the image of the $G$--stabilizer of $[f]$ in $\mathrm{PGL}_n$ is finite; its kernel
is contained in the scalar matrices of $G$, the finite group $\mu_n$, so the stabilizer
itself is finite.

\emph{Stability.} A smooth hypersurface of degree $d\ge 3$ is a stable point for the
$G$--action on $\PP(\mathcal V)$ in the sense of geometric invariant theory
\cite[\S 4.2]{Mumford}, \cite[Ch.~1, \S 4]{MFK}. A stable point with finite stabilizer is
properly stable. Thus $\PP X$ lies in the properly stable locus $\PP(\mathcal V)^{s}$.

\emph{Geometric quotient and normality.} By the fundamental theorem of GIT
\cite[Ch.~1, Thm.~1.10 and Converse~1.13]{MFK}, the restriction of the quotient map to the
properly stable locus is a geometric quotient onto an open subvariety of
$\PP(\mathcal V)/\!/G$, and $\PP X\to M$ is its restriction over the open $G$--invariant
$\PP X$. So $M$ is quasi-projective and $\PP X\to M$ is a geometric quotient. Finally $\PP X$
is smooth (open in the projective space $\PP(\mathcal V)$, complement of the discriminant),
hence normal; a geometric quotient of a normal variety by a reductive group is normal
\cite[Ch.~1, \S 2]{MFK}, so $M$ is normal.
\end{proof}

\section{Euler's identity splits the gradient}\label{sec:euler}

We record the equivariance of $\nab$ and its splitting. Define the $\GL_n$--action on
$\mathcal W=(A_{d-1})^n$ by
\begin{equation}\label{eq:Zaction}
g\cdot (p_1,\dots,p_n)
:= \Big(\textstyle\sum_{j}(g^{-T})_{1j}\,(g\cdot p_j),\ \dots,\ \sum_{j}(g^{-T})_{nj}\,(g\cdot p_j)\Big),
\quad g^{-T}=(g^{-1})^{T},
\end{equation}
$g\cdot p_j$ being the form $v\mapsto p_j(g^{-1}v)$. This is the natural action making the
gradient equivariant and is the one carried by $Z$ and $N$.

\begin{lemma}\label{lem:nabeq}
$\nab\colon \mathcal V\to\mathcal W$ is $\GL_n$--equivariant: $\nab(g\cdot f)=g\cdot\nab f$
for all $g\in\GL_n$, $f\in\mathcal V$.
\end{lemma}

\begin{proof}
With $g\cdot f=f\circ g^{-1}$, the chain rule gives for each $i$
\[
\partial_i(g\cdot f)(v)=\sum_j (g^{-1})_{ji}\,(\partial_j f)(g^{-1}v)
=\sum_j (g^{-T})_{ij}\,(g\cdot\partial_j f)(v),
\]
which is the $i$-th component of $g\cdot\nab f$ in \eqref{eq:Zaction}.
\end{proof}

\begin{lemma}\label{lem:euler}
The linear map
\[
R\colon \mathcal W\to\mathcal V,\qquad R(p_1,\dots,p_n)=\frac1d\sum_{i=1}^n x_i\,p_i,
\]
satisfies:
\textup{(a)} $R\circ\nab=\mathrm{id}_{\mathcal V}$;
\textup{(b)} $R$ is $\GL_n$--equivariant for the actions \eqref{eq:Zaction} on $\mathcal W$
and $g\cdot h=h\circ g^{-1}$ on $\mathcal V$.
\end{lemma}

\begin{proof}
(a) Euler's identity for a degree-$d$ form: $\sum_i x_i\partial_i f=d\,f$, so
$R(\nab f)=\tfrac1d\sum_i x_i\partial_i f=f$.

(b) With $g\cdot p=(q_1,\dots,q_n)$, $q_i=\sum_j(g^{-T})_{ij}(g\cdot p_j)$,
\[
d\,R(g\cdot p)=\sum_i x_i q_i=\sum_j (g\cdot p_j)\sum_i (g^{-1})_{ji}x_i
=\sum_j (g\cdot p_j)\,(g^{-1}x)_j .
\]
Evaluating at $v$: $(g\cdot p_j)(v)=p_j(g^{-1}v)$ and $(g^{-1}x)_j(v)=(g^{-1}v)_j$, so the
right side equals $\big(\sum_j x_j p_j\big)(g^{-1}v)=d\,(g\cdot R(p))(v)$. Hence
$R(g\cdot p)=g\cdot R(p)$.
\end{proof}

Thus $\nab$ is a $\GL_n$--equivariant linear injection that is \emph{split} by the
$\GL_n$--equivariant linear map $R$.

\section{Surjectivity of the pullback on invariant rings}\label{sec:inv}

Let $\mathcal V^{*},\mathcal W^{*}$ be the dual spaces (linear coordinate functions), and let
\[
\CC[\mathcal V]=\Sym^{\bullet}(\mathcal V^{*}),\qquad \CC[\mathcal W]=\Sym^{\bullet}(\mathcal W^{*})
\]
be the polynomial coordinate rings of the affine spaces $\mathcal V,\mathcal W$. We grade them
by ordinary polynomial degree (this is the grading by the overall scaling $\CC^{*}$ on
$\mathcal V$ resp.\ $\mathcal W$). The linear map $\nab\colon \mathcal V\to\mathcal W$ induces
the graded $\CC$--algebra pullback homomorphism
\[
\nab^{\#}\colon \CC[\mathcal W]\to\CC[\mathcal V],\qquad \nab^{\#}(F)=F\circ\nab,
\]
which is $G$--equivariant by Lemma~\ref{lem:nabeq} (where $G$ acts on coordinate rings via the
contragredient of its action on $\mathcal V,\mathcal W$).

\begin{lemma}\label{lem:nabsharp-surj}
$\nab^{\#}\colon \CC[\mathcal W]\to\CC[\mathcal V]$ is surjective.
\end{lemma}

\begin{proof}
It suffices to hit the degree-one part $\mathcal V^{*}$, since $\CC[\mathcal V]$ is generated
as a $\CC$--algebra by $\mathcal V^{*}$ and $\nab^{\#}$ is a graded ring homomorphism.
Let $\ell\in\mathcal V^{*}$ be any linear functional on $\mathcal V$. Consider the linear
functional $F:=\ell\circ R\in\mathcal W^{*}$, where $R\colon\mathcal W\to\mathcal V$ is the
Euler splitting (Lemma~\ref{lem:euler}). Then $F\in\CC[\mathcal W]$ has degree one and, by
Lemma~\ref{lem:euler}(a),
\[
\nab^{\#}(F)=F\circ\nab=\ell\circ R\circ\nab=\ell\circ\mathrm{id}_{\mathcal V}=\ell .
\]
Hence $\mathcal V^{*}\subset\im\nab^{\#}$, and $\nab^{\#}$ is surjective.
\end{proof}

\begin{lemma}\label{lem:phi-surj}
The restriction of $\nab^{\#}$ to $G$--invariants,
\[
\varphi:=\nab^{\#}|_{\CC[\mathcal W]^{G}}\colon \CC[\mathcal W]^{G}\longrightarrow
\CC[\mathcal V]^{G},
\]
is a well-defined surjective graded $\CC$--algebra homomorphism.
\end{lemma}

\begin{proof}
Since $\nab^{\#}$ is $G$--equivariant it maps invariants to invariants, so $\varphi$ is
well defined; it is a graded algebra homomorphism because $\nab^{\#}$ is. For surjectivity,
let $K=\ker\nab^{\#}\subset\CC[\mathcal W]$, a $G$--stable graded ideal. By
Lemma~\ref{lem:nabsharp-surj} there is a short exact sequence of graded
$G$--representations
\[
0\to K\to \CC[\mathcal W]\xrightarrow{\ \nab^{\#}\ }\CC[\mathcal V]\to 0 .
\]
Over $\CC$ the group $G=\SL_n$ is linearly reductive, so the functor $(-)^{G}$ of taking
$G$--invariants is exact on the category of $G$--representations
\cite[Ch.~1, \S 2, Reductivity]{MFK}. Applying it yields the exact sequence
\[
0\to K^{G}\to \CC[\mathcal W]^{G}\xrightarrow{\ \varphi\ }\CC[\mathcal V]^{G}\to 0 ,
\]
in particular $\varphi$ is surjective.
\end{proof}

\section{From surjective invariant pullback to an isomorphism onto the image}\label{sec:proj}

Write $S:=\CC[\mathcal W]^{G}$ and $R^{\bullet}:=\CC[\mathcal V]^{G}$ for the two graded
invariant rings; both are finitely generated $\CC$--algebras with $S_0=R^{\bullet}_0=\CC$
(Hilbert--Nagata for the reductive group $G$ \cite[Ch.~1, \S 2]{MFK}). By construction
\[
\PP(\mathcal V)/\!/G=\Proj R^{\bullet},\qquad \PP(\mathcal W)/\!/G=\Proj S,
\]
and the GIT quotient $N=\PP Z/\!/G$ is the open subscheme of $\Proj S$ over the semistable
locus $\PP Z$, while $\PP(\mathcal V)/\!/G=\Proj R^{\bullet}$ contains $M$ as the open
subvariety corresponding to the stable locus $\PP X$ (Proposition~\ref{prop:sourcenormal}).

\begin{lemma}\label{lem:closedimm}
The graded surjection $\varphi\colon S\twoheadrightarrow R^{\bullet}$ of
Lemma~\ref{lem:phi-surj} induces a closed immersion of projective schemes
\[
\iota\colon \Proj R^{\bullet}\hookrightarrow \Proj S ,
\]
and on the affine level a closed immersion of affine GIT quotients
$\Spec R^{\bullet}\hookrightarrow\Spec S$.
\end{lemma}

\begin{proof}
A surjective homomorphism of finitely generated graded $\CC$--algebras
$\varphi\colon S\twoheadrightarrow R^{\bullet}$ (with $\varphi$ preserving the grading and
$S_0=R^{\bullet}_0=\CC$) induces a morphism $\iota\colon\Proj R^{\bullet}\to\Proj S$ defined
on all of $\Proj R^{\bullet}$: for a homogeneous prime $\mathfrak p\subset R^{\bullet}$ not
containing the irrelevant ideal $R^{\bullet}_{+}$, the preimage $\varphi^{-1}(\mathfrak p)$ is
a homogeneous prime of $S$ not containing $S_{+}$ (because $\varphi(S_{+})=R^{\bullet}_{+}$ by
surjectivity, so $S_{+}\subset\varphi^{-1}(\mathfrak p)$ would force
$R^{\bullet}_{+}\subset\mathfrak p$, a contradiction). On each basic open
$D_{+}(\varphi(s))\subset\Proj R^{\bullet}$, $s\in S_{+}$ homogeneous, $\iota$ is the morphism
$\Spec(R^{\bullet}_{(\varphi(s))})\to\Spec(S_{(s)})$ dual to the ring map
$S_{(s)}\to R^{\bullet}_{(\varphi(s))}$ induced by $\varphi$, which is surjective (localizing a
surjection). A morphism that is, locally on the target, $\Spec$ of a surjective ring map is a
closed immersion; covering $\Proj S$ by the $D_{+}(s)$ and noting $\iota^{-1}(D_{+}(s))\subset
D_{+}(\varphi(s))$, we conclude $\iota$ is a closed immersion. The affine statement
$\Spec R^{\bullet}\hookrightarrow\Spec S$ is immediate from surjectivity of $\varphi$.
\end{proof}

\begin{proposition}\label{prop:iso-onto-image}
$\PP\nab\colon M\to N$ is an isomorphism onto its image $Y$; equivalently $\PP\nab$ is a
locally closed immersion with closed image.
\end{proposition}

\begin{proof}
By Lemma~\ref{lem:nabeq} the linear map $\nab$ is $G$--equivariant and homogeneous of degree
$1$ for the scalings, so it induces the morphism of GIT quotients
$\Phi\colon \Proj R^{\bullet}=\PP(\mathcal V)/\!/G\to \PP(\mathcal W)/\!/G=\Proj S$ that is the
$\Proj$ of $\varphi$; this is exactly the map $\iota$ of Lemma~\ref{lem:closedimm}. Indeed
the morphism induced on GIT quotients by a $G$--equivariant graded-degree-one linear map is,
by definition of the GIT quotient as $\Proj$ of the invariant ring, the $\Proj$ of the
induced map on invariant rings $\varphi=\nab^{\#}|_{(-)^G}$. By Lemma~\ref{lem:closedimm},
$\Phi=\iota$ is a closed immersion of $\Proj R^{\bullet}$ into $\Proj S$.

Now restrict to the relevant open loci. The map $\PP\nab$ is the restriction of $\Phi$ to the
open subvariety $M\subset \Proj R^{\bullet}$ (the stable locus), landing in the open
subvariety $N\subset\Proj S$ (the semistable locus): this is precisely the compatibility
``$\PP\nab$ is induced by $\nab$'' and the fact that $\nab$ carries (semi)stable points of
$\mathcal V$ to (semi)stable points of $\mathcal W$, which is built into the hypothesis that
$\PP\nab\colon M\to N$ is defined. A closed immersion restricted to an open subscheme of the
source and corestricted to an open subscheme of the target is a locally closed immersion;
hence $\PP\nab\colon M\to N$ is a locally closed immersion, i.e.\ an isomorphism of $M$ onto a
locally closed subvariety $\PP\nab(M)$ of $N$.

Finally $\PP\nab$ is finite (given), hence proper, hence closed; so its image $Y=\PP\nab(M)$
is closed in $N$. A locally closed immersion with closed image is a closed immersion, and in
particular $\PP\nab\colon M\to Y$ is an isomorphism.
\end{proof}

\section{Proof of the theorem}\label{sec:conclude}

\begin{proof}[Proof of Theorem~\ref{thm:main}]
By Proposition~\ref{prop:iso-onto-image}, $\PP\nab\colon M\to Y$ is an isomorphism onto its
closed image $Y\subset N$. By Proposition~\ref{prop:sourcenormal}, $M$ is normal. Normality
is preserved under isomorphism, so $Y\cong M$ is normal. This proves Theorem~\ref{thm:main}
for all $n\ge 2$ and $d\ge 3$.
\end{proof}

\begin{remark}[Consistency with the problem's remark]
The problem notes that normality of $Y$ would, via Zariski's Main Theorem, imply $\PP\nab$ is
an isomorphism onto its image. We verify there is no circularity: we did \emph{not} assume
normality of $Y$ to obtain the isomorphism. The isomorphism in
Proposition~\ref{prop:iso-onto-image} is derived independently, from the algebraic fact that
$\nab$ admits a $G$--equivariant linear left inverse $R$ (Euler's identity), which forces
$\nab^{\#}$ and hence (by linear reductivity of $G$) the induced map on invariant rings
$\varphi$ to be surjective; surjectivity of $\varphi$ yields a closed immersion of GIT
quotients. Normality of $Y$ is then a consequence, not a hypothesis. Conversely, given
normality, Zariski's Main Theorem reproves the isomorphism: $\PP\nab$ is finite, injective,
and birational onto $Y$ (a finite injective morphism in characteristic $0$ is birational onto
its image), and a finite birational morphism onto a normal variety is an isomorphism
\cite[Ch.~III, \S 9]{Mumford-RB}.
\end{remark}

\begin{thebibliography}{9}

\bibitem{MM} H. Matsumura and P. Monsky,
\emph{On the automorphisms of hypersurfaces}, J. Math. Kyoto Univ. \textbf{3} (1963/64),
347--361.

\bibitem{Mumford} D. Mumford,
\emph{Stability of projective varieties}, Enseign. Math. (2) \textbf{23} (1977), 39--110.

\bibitem{MFK} D. Mumford, J. Fogarty, F. Kirwan,
\emph{Geometric Invariant Theory}, 3rd ed., Ergeb. Math. Grenzgeb. (2) \textbf{34},
Springer, Berlin, 1994.

\bibitem{Mumford-RB} D. Mumford,
\emph{The Red Book of Varieties and Schemes}, 2nd ed., Lecture Notes in Math. \textbf{1358},
Springer, Berlin, 1999.

\end{thebibliography}

\end{document}
